% once the ethical issue(s) is/are identified and the impact and stakeholders are known, one or more design directions are considered to making the AI system responsible: value alignment, Explainable AI, governance etc.


% \subsection{Design Introduction}
% introduce the type of method you plan to consider applying corresponding controls and transform the existing system into a responsible AI system.


% \subsection{Design Description}
% discuss the responsible AI methods applied using one or more of the instruments of choice (e.g., just text, design architecture, short implementation etc.)

De literatuur met `NEW' erbij staan reeds in de bronnen van de case (in het bestand RAI Assignment 2 Description) maar hebben we nog niet besproken omdat we nog niet bezig waren met solutions

\subsection{Technical layer (XAI, interpretability)}
\begin{itemize}
	\item \textbf{\citet{arrieta2019explainableartificialintelligencexai}}: post-hoc solutions (SHAP, LIME, attention visualisation) applied to ensemble models
	\item \textbf{\citet{miller2018explanationartificialintelligenceinsights}}: XAI
	\item \textbf{(NEW) Liao, Gruen \& Miller (2020)~\citep{liao2020questioning}}: XAI question bank (Why this target? What evidence supports this recommendation?) $\rightarrow$ design requirements
	\item \textbf{(NEW) Morley et al. (2020)~\citep{morley2020what}}: maps ethics tools onto ML pipeline stages (data, model, deployment) so we can use this to say where and how speicifc XAI tools can be added before/during/after model development
\end{itemize}


\subsection{Socio-technical layer (MHC, automation bias, VSD, stakeholder inclusion)}\
\begin{itemize}
	\item \textbf{\citet{mhcof}}: The 2 MHC conditions (sufficient information, meaningful freedom of action) $\rightarrow$ design requirements
	\item \textbf{\citet{parasuraman2010complacency}}: VSD's `tripartite' methodology (conceptual, empirical, technical) $\rightarrow$ concrete process for how stakeholders shold have been consulted (Palestinians were not consulted)
	% \item \textbf{(NEW) Friedman, Kahn, Borning \& Huldtgren (2013)~\citep{friedman2013vsd}}: % TODO
	\item \textbf{(NEW) Amershi et al. (2019)~\citep{amershi2019guidelines}}: human-AI interaction guidelines (calibration of trust, commuinicating uncertainty, supporting error recovery) $\rightarrow$ design principles
\end{itemize}


\subsection{Governance (accountability, legal frameworks, organisational response)}
\begin{itemize}
	\item \textbf{\citet{info12090385}}: addresses the empty cell in Figure 1 $\rightarrow$ required governance interventions (pre-deployment EIA, real-time IHL compliance, post-deployment audit)
	\item \textbf{(NEW) Stahl et al. (2022)~\citep{stahl2022organisational}}: Use the 3-level mitigation (policy-level strategies, guidance mechanisms, corporate governance) to categorise governance solutions
	\item \textbf{\citet{aimag.v40i1.2848}}: Canada has pre-deployment legal assessment protocols (contrast with Israel that does not, because it does not follow Article 36 review process)
	\item \textbf{\citet{khan2021ethicsaisystematicliterature}}: P-01/P-03/CH-06/CH-09 $\rightarrow$ we can use this to motiviate specific governance solutions
\end{itemize}
